
\chapter{جمع‌بندی و کارهای آتی}
\section{جمع‌بندی}
در این گزارش تلاش شد تا اطلاعات اولیه‌ای برای آشنایی و یافتن یک دیدگاه کلی در خصوص مدل‌های پخشی ارائه شود. هدف کلی از این گزارش بررسی مقالات پایه‌ و مبنای تئوری این مدل‌ها برای درک ساختار و سازو کار تولید داده بود و از پرداختن به حجم زیادی از کاربردها و بهبودها که با توجه به رشد روز افزون به کارگیری مدل های پخشی رو به فزونی است خودداری شده است.
\\
در فصل اول به معرفی مدل‌های پخشی مبتنی بر امتیاز پرداختیم. این مدل‌ها در رویکرد اولیه‌ی خود تماما پخشی نیستند و تنها با استفاده از یک تطبیق امتیاز و دینامیک لانژوین قابل انجام هستند. اما رویکرد اولیه‌ی این روش‌ها با محدودیت‌هایی از جمله مشکل رویه‌ها و نواحی کمتر چگال توزیع داده روبرو بود که این محدودیت‌ها در داده‌های واقعی اکثر اوقات برقرار هستند. به همین دلیل در مدل ارائه شده که تطبیق امتیاز نویززدا با دینامیک لانژوین تابکاری شده نام دارد، به دنبال برطرف کردن این مشکل هستیم و این روش مانند تمام رویکرد‌های پخشی به صورت مرحله به مرحله داده را با نویز تخریب و سپس در فرآیند معکوس نمونه‌هایی از داده را تولید می‌کند.
\\
در فصل دوم به معرفی مپنا پرداختیم، مدلی که رایج‌ترین نوع مدل‌های پخشی است و در کاربرد‌های تولید تصویر مطرح مانند $Dall-E$ به کار گرفته شده است. این مدل به صورت احتمالاتی و با یک زنجیره ی مارکفی فرآیند پخش را مدل می‌کند. برای آموزش این شبکه رویکرد های مختلفی وجود دارد. در یک رویکرد می‌توانیم از تخمین میانگین پسین فرآیند رو به جلو برای آموزش شبکه استفاده کنیم و در رویکرد دوم می‌توانیم مستقیما نویز اضافه شده به تصویر در یک گام از فرآیند پخش را تخمین بزنیم. می‌توان نشان داد که با تغییر پارامتر در رویکرد دوم و ثابت در نظر گرفتن واریانس این روش هم ارز روش تطبیق امتیاز با دینامیک لانژوین است.
\\
در فصل سوم چهارچوبی کلی برای فرآیند پخشی معرفی شد که قابلیت دربرگیری روش‌های پیش‌تر معرفی شده را دارد. در این روش نیز به صورت مرحله به مرحله داده را با نویز تخریب می‌کنیم اما به جای به کارگیری تعداد گام‌های محدود از تعداد گام نامتناهی پیوسته در یک بازه استفاده می‌کنیم. برای مدل کردن فرآیند پخش در چنین حالتی از یک معادله دیفرانسیل تصادفی استفاده می‌کنیم که معکوس چنین معادله‌ای را می‌توانیم تنها با داشتن تابع امتیاز بسازیم. برای ساخت نمونه از چنین معادله‌ای تنها کافی است از یک حل کننده استفاده کنیم و آن را گسسته کنیم.  همچنین در این فصل نشان دادیم که هر دو روش پیش‌تر معرفی شده با تعیین معادلات متناسب با آنها می‌توانند در چهارچوب روش‌های مبتنی بر معادله دیفرانسیل تصادفی قرار گیرند.
\\
در فصل چهارم پس از معرفی سه دسته‌ی اصلی مدل‌های پخشی به مقایسه‌ی آنها هم با سایر مدل‌های مولد و هم با یکدیگر پرداختیم. نکته‌ی قابل تاکید در این فصل این است که برای مقایسه‌ی روش‌های معرفی شده تنها به مقالات به کار برده شده در سراسر گزارش اتکا شده است و با توجه به استقبال روزافزون از توسعه و به کارگیری این روش‌ها ممکن است به نتایج و بهبود‌های بهترو یا رفع مشکلات مطرح شده پرداخته شده باشد که شرح آنها در این گزارش نیامده است.
\\
\section{کارهای آتی}
با توجه به مطالعات انجام شده و با توجه به نوظهور بودن این دسته از مدل‌ها قابلیت بهبود و توسعه برای آنها در انواعی از جنبه ها وجود دارد که در ادامه آنها را نام می‌بریم.

\begin{itemize}
	\item
	\textbf{تئوری}
	\\
	از آنجاییکه مدل‌های پخشی اساسا بر پایه‌ی روابط ریاضی توسعه داده شده اند، امکان بهبود آنها از جنبه ی ریاضی و تئوری وجود دارد. به عنوان مثال ارائه‌ی فرمولاسیون جدیدی برای تخریب داده‌ها، استفاده از انواعی از حل‌کننده‌های معادله دیفرانسیل و بررسی تاثیرات آنها و همچنین به کارگیری مباحثی از فیزیک ترمودینامیک از جمله روش‌های قابل اعمال هستند.
	\item
	\textbf{معماری}
	\\
	معماری فعلی تمام مدل‌های پخشی از هسته‌ی شبکه‌ی یو بهره می‌برد. با وجود عملکرد فوق‌العاده‌ی این شبکه‌ها همچنان مشکلاتی در استفاده از آنها در بعضی از کاربردهای خاص وجود دارد و علاوه بر این تعداد پارامترهای این شبکه‌ها آنها را از لحاظ محاسباتی بسیار سنگین می‌کند. پیاده سازی و استفاده از سایر شبکه‌ها برای شبکه‌ی نویز زدا می‌تواند پیشرفت قابل توجهی در به کارگیری این مدل‌ها در همه‌ی کاربردها فراهم آورد.
	\item
	\textbf{داده}
	\\
	آموزش و توسعه‌ی مدل‌های پخشی با میزان داده‌ی کارا یک چالش باز در بحث توسعه‌ی مدل‌های پخشی است. این مدل‌ها برای تولید داده‌ی با کیفیت نیاز به دریافت داده‌ی با کیفیت دارند و از این رو کار روی یادگیری محدود و یادگیری با داده‌ی کم کیفیت می تواند یک رویکرد قابل پیگیری باشد.
	\item
	\textbf{کاربرد}
	\\
	همانطور که دراین گزارش نیز مشاهده شد، اکثر توسعه‌های مدل‌های پخشی بر روی داده های تصویری انجام شده است. اما تحقیقات متعددی روی پیاده سازی و تطبیق این مدل‌ها در سایرحوزه‌های کاربرد مانند صوت، متن، یادگیری تقویتی و گراف نیز در حال انجام است. تطبیق مدل‌های پخشی با انواع اشکال داده گونه‌ای دیگر از مسائلی است که می‌تواند مورد بررسی و تغییر قرار گیرد. 
